% ============================================================
%  ADD-ON CONTENT FOR GeoSI_Master_Report.tex
%  ── Place this chapter BEFORE "Conclusions and Future Work"
%  ── i.e., as Chapter 10 (renumber Conclusions to Chapter 11)
% ============================================================


% ============================================================
%  CHAPTER 10 — NATIVE HYDROLOGICAL ANALYSIS
% ============================================================
\chapter{Native Hydrological Analysis Pipeline}

\section{Overview}

A key extension to GeoSI is the implementation of a dependency-free, native
hydrological analysis pipeline within the \texttt{WatershedWorkflowTool}
(\texttt{geosi\_engine/tools/complex.py}).  Traditional hydrological analysis in
QGIS requires external providers such as GRASS GIS or SAGA, which are often
difficult to install and not available on all platforms.  GeoSI replaces these
external dependencies with a pure NumPy and GDAL implementation of the D8 flow
direction algorithm, enabling watershed delineation and pour-point catchment
analysis to run natively on any platform with only \texttt{osgeo.gdal} and
\texttt{numpy}.

\section{D8 Flow Direction Algorithm}

The D8 (Deterministic Eight-Neighbour) algorithm assigns a single flow direction
to each cell in a Digital Elevation Model (DEM) based on the steepest downhill
slope among its eight neighbours.

\subsection{Mathematical Formulation}

For a cell at position $(r, c)$ with elevation $z_{r,c}$, the slope to each
neighbour $(r + \Delta r,\; c + \Delta c)$ is computed as:

\[
S_{(\Delta r, \Delta c)} = \frac{z_{r,c} - z_{r+\Delta r,\; c+\Delta c}}{d_{(\Delta r,\Delta c)}}
\]

\noindent where $d$ is the distance to the neighbour: $d = 1$ for cardinal
directions (N, S, E, W) and $d = \sqrt{2}$ for diagonal directions (NE, SE, SW,
NW).  The flow direction is assigned as the D8 code of the neighbour with the
maximum positive slope.

\subsection{D8 Direction Encoding}

GeoSI uses the standard D8 encoding:

\begin{table}[ht]
\centering
\caption{D8 flow direction codes}
\begin{tabular}{cccccccc}
\toprule
\textbf{E} & \textbf{SE} & \textbf{S} & \textbf{SW} & \textbf{W} & \textbf{NW} & \textbf{N} & \textbf{NE} \\
\midrule
1 & 2 & 4 & 8 & 16 & 32 & 64 & 128 \\
\bottomrule
\end{tabular}
\end{table}

\subsection{Vectorised Implementation}

Rather than looping over every pixel in Python (which would be prohibitively slow
for large DEMs), GeoSI uses a vectorised sliding-window approach with NumPy:

\begin{lstlisting}
# Compute slope to all 8 neighbours simultaneously
for i, (dc, dr, code) in enumerate(D8_OFFSETS):
    neighbor = dem_padded[1+dr:rows+1+dr, 1+dc:cols+1+dc]
    dist = math.sqrt(2) if abs(dc)+abs(dr)==2 else 1.0
    slope = (center - neighbor) / (dist * pixel_size)
    slopes[:, :, i] = slope
    codes[i] = code

# Assign direction = code of steepest downhill neighbour
best = np.argmax(slopes, axis=2)
flow_dir = np.array(codes)[best]
flow_dir[np.max(slopes, axis=2) <= 0] = 0  # pits / flats
\end{lstlisting}

This processes the entire DEM in a single pass with no Python-level pixel loops,
achieving near-C performance on large rasters.

\section{Flow Accumulation}

Flow accumulation counts the number of upstream cells that drain through each
cell.  GeoSI computes this using a topological sort by elevation:

\begin{enumerate}
  \item Sort all pixels in descending elevation order.
  \item Initialise an accumulation grid with value 1 for every cell.
  \item Iterate from highest to lowest: for each cell, add its accumulation
        value to the cell it flows into (determined by its D8 direction).
\end{enumerate}

This avoids the deep recursion that would cause Python's stack to overflow on
large DEMs.  The topological sort runs in $O(n \log n)$ where $n$ is the total
number of pixels, and the accumulation pass runs in $O(n)$.

\begin{lstlisting}
# Topological sort: process cells from highest to lowest
sorted_indices = np.argsort(-dem_filled.ravel())
accumulation = np.ones(dem_filled.shape, dtype=np.float64)

for idx in sorted_indices:
    r, c = divmod(idx, cols)
    direction = flow_dir[r, c]
    if direction == 0:
        continue
    dr, dc = D8_DECODE[direction]
    nr, nc = r + dr, c + dc
    if 0 <= nr < rows and 0 <= nc < cols:
        accumulation[nr, nc] += accumulation[r, c]
\end{lstlisting}

\section{Watershed Delineation}

Watersheds are delineated by thresholding the flow accumulation grid.  Cells
with accumulation above the user-specified threshold are marked as part of the
drainage network, then converted to vector polygons using GDAL's
\texttt{Polygonize} function.

\begin{lstlisting}
# Threshold to create binary watershed mask
mask = (accumulation >= threshold).astype(np.uint8)

# Vectorise using GDAL
gdal.Polygonize(mask_band, mask_band, output_layer, 0)
\end{lstlisting}

\section{Pour-Point Catchment Delineation}

The most significant hydrological feature is pour-point catchment delineation:
given a specific coordinate (the ``pour point''), the tool identifies all upstream
cells that drain into that point.

\subsection{Pour Point Snapping}

Users rarely click exactly on the mathematical stream channel.  The algorithm
first snaps the input coordinate to the cell with the highest flow accumulation
within a configurable search radius:

\begin{lstlisting}
# Search window around the pour point
sub_fa = fa_data[min_y:max_y, min_x:max_x]
max_idx = np.unravel_index(np.argmax(sub_fa), sub_fa.shape)
snap_py = min_y + max_idx[0]
snap_px = min_x + max_idx[1]
\end{lstlisting}

\subsection{Reverse D8 Upstream Trace}

Starting from the snapped pour point, the algorithm performs an iterative
depth-first search (DFS) upstream.  For each of the eight neighbours, it checks
whether that neighbour's D8 flow direction points \emph{towards} the current
cell.  If so, the neighbour is part of the catchment:

\begin{lstlisting}
# Reverse D8 neighbour lookup
# If neighbour at (dx, dy) has flow code = required_code,
# it flows INTO the current cell
neighbours = [
    (0, -1, 4),    # N neighbour must flow S (code 4)
    (1, -1, 8),    # NE neighbour must flow SW (code 8)
    (1,  0, 16),   # E neighbour must flow W (code 16)
    (1,  1, 32),   # SE neighbour must flow NW (code 32)
    (0,  1, 64),   # S neighbour must flow N (code 64)
    (-1, 1, 128),  # SW neighbour must flow NE (code 128)
    (-1, 0, 1),    # W neighbour must flow E (code 1)
    (-1,-1, 2)     # NW neighbour must flow SE (code 2)
]

stack = [(snap_px, snap_py)]
catchment_mask[snap_py, snap_px] = 1

while stack:
    cx, cy = stack.pop()
    for dx, dy, required_code in neighbours:
        nx, ny = cx + dx, cy + dy
        if 0 <= nx < cols and 0 <= ny < rows:
            if catchment_mask[ny, nx] == 0:
                if flow_dir[ny, nx] == required_code:
                    catchment_mask[ny, nx] = 1
                    stack.append((nx, ny))
\end{lstlisting}

The use of an explicit stack (rather than recursion) prevents stack overflow on
large DEMs where catchments may contain millions of pixels.

\subsection{Multiple Pour Points from Shapefile}

The tool supports reading pour points from a vector layer (e.g., a point
shapefile) using OGR.  When a \texttt{POUR\_POINT\_LAYER} parameter is provided,
the tool reads all point geometries from the file and runs a separate catchment
delineation for each point:

\begin{lstlisting}
def _resolve_pour_points(self, kwargs):
    points = []
    pp_layer = kwargs.get("POUR_POINT_LAYER")
    if pp_layer is not None:
        ds = ogr.Open(pp_layer.filepath)
        layer = ds.GetLayer(0)
        for feature in layer:
            geom = feature.GetGeometryRef()
            points.append((geom.GetX(), geom.GetY()))
    # Fallback: explicit X/Y coordinates
    if not points:
        pt_x = kwargs.get("POUR_POINT_X")
        pt_y = kwargs.get("POUR_POINT_Y")
        if pt_x is not None and pt_y is not None:
            points.append((float(pt_x), float(pt_y)))
    return points
\end{lstlisting}

\section{Unified Workflow Pipeline}

The \texttt{WatershedWorkflowTool} orchestrates the entire pipeline as a single
tool invocation with up to six steps:

\begin{table}[ht]
\centering
\caption{WatershedWorkflowTool execution pipeline}
\begin{tabular}{clp{7cm}}
\toprule
\textbf{Step} & \textbf{Operation} & \textbf{Output} \\
\midrule
1 & Preprocess DEM        & Filled DEM with no-data values interpolated \\
2 & D8 Flow Direction     & 8-bit raster with flow codes (1--128) \\
3 & Flow Accumulation     & Float64 raster with upstream cell counts \\
4 & Watershed Mask        & Binary raster (above/below threshold) \\
5 & Polygonise            & GeoJSON vector of watershed boundaries \\
6 & Catchment (optional)  & Catchment mask + polygon for each pour point \\
\bottomrule
\end{tabular}
\end{table}

Steps 1--5 always execute.  Step 6 runs only when pour point coordinates or a
pour point layer are provided.

\section{Natural Language Integration}

The parser and agent were updated to support watershed queries from the GeoSI
plugin chat panel:

\subsection{Parser Updates}

The \texttt{\_OPERATION\_MAP} in \texttt{parser.py} was extended with a new
\texttt{"watershed"} entry:

\begin{lstlisting}
"watershed": ["watershed", "watershed analysis",
              "flow accumulation", "flow direction",
              "catchment", "drainage basin",
              "pour point", "delineate"],
\end{lstlisting}

A coordinate extraction regex was added to parse pour point coordinates from
natural language:

\begin{lstlisting}
coord_match = re.search(
    r'(?:at\s+(?:point\s+)?|pour\s+point\s+|'
    r'coordinates?\s+|point\s+)'
    r'(-?\d+\.?\d*)[,\s]+(-?\d+\.?\d*)',
    lowered
)
\end{lstlisting}

\subsection{Agent Updates}

The agent's \texttt{op\_to\_tool} mapping was extended:

\begin{lstlisting}
"watershed": "watershed_workflow",
\end{lstlisting}

When the parser detects pour point coordinates or a secondary layer (point
shapefile), the agent passes them as \texttt{POUR\_POINT\_X}/\texttt{Y} or
\texttt{POUR\_POINT\_LAYER} parameters.

\subsection{Supported Prompts}

\begin{table}[ht]
\centering
\caption{Supported natural language prompts for hydrological analysis}
\begin{tabular}{lp{6cm}}
\toprule
\textbf{Prompt} & \textbf{Behaviour} \\
\midrule
``Perform watershed analysis on DEM'' & Full 5-step watershed (no catchment) \\
``Watershed analysis on DEM at point 77.5, 8.5'' & Watershed + catchment at (77.5, 8.5) \\
``Watershed analysis on DEM with points'' & Watershed + catchment for each point in the \texttt{points} layer \\
``Find catchment at 77.5, 8.5 on DEM'' & Watershed + catchment at (77.5, 8.5) \\
\bottomrule
\end{tabular}
\end{table}

\section{Case Study: Watershed and Catchment on Kerala DEM}

The pipeline was tested on a 3601$\times$3601 pixel DEM (SRTM 1-arc-second)
covering Kerala (bounds: 76.99--78.00\textdegree E, 8.00--9.00\textdegree N).

\subsection{Execution Results}

\begin{table}[ht]
\centering
\caption{Watershed workflow execution on Kerala DEM (threshold=5000)}
\begin{tabular}{lcr}
\toprule
\textbf{Output Layer} & \textbf{Type} & \textbf{Size} \\
\midrule
\texttt{dem\_filled}            & GeoTIFF & 25.9\,MB \\
\texttt{flow\_direction}        & GeoTIFF & 12.9\,MB \\
\texttt{flow\_accumulation}     & GeoTIFF & 51.8\,MB \\
\texttt{watershed\_mask}        & GeoTIFF & 12.9\,MB \\
\texttt{watershed\_polygons}    & GeoJSON & 185\,B \\
\texttt{catchment\_mask}        & GeoTIFF & 12.9\,MB \\
\texttt{catchment\_polygon}     & GeoJSON & 5.7\,KB \\
\bottomrule
\end{tabular}
\end{table}

All seven layers were successfully generated and loaded into the QGIS Layers
panel via the GeoSI plugin in approximately 25 seconds total execution time.

\subsection{QGIS Plugin Execution}

The watershed workflow was triggered from the GeoSI plugin chat panel using the
prompt: ``Perform watershed analysis on DEM with a.shp'' (where \texttt{a.shp}
contained a single pour point).  The plugin successfully:

\begin{enumerate}
  \item Parsed the query as \texttt{intent=TERRAIN}, \texttt{operation=watershed}.
  \item Resolved \texttt{DEM} as the primary layer and \texttt{a.shp} as the pour
        point layer.
  \item Executed all six pipeline steps using the native D8 algorithm.
  \item Loaded all seven output layers into the QGIS Layers panel automatically.
\end{enumerate}

\section{Design Decisions}

\begin{enumerate}
  \item \textbf{No external dependencies.}  The pipeline uses only
        \texttt{osgeo.gdal} (bundled with QGIS) and \texttt{numpy}.  No GRASS,
        SAGA, or WhiteboxTools installation is required.
  \item \textbf{Iterative stack, not recursion.}  The upstream trace uses an
        explicit Python list as a stack to prevent \texttt{RecursionError} on
        large catchments.
  \item \textbf{Topological sort for accumulation.}  Sorting by descending
        elevation guarantees correct accumulation without recursive descent.
  \item \textbf{Multi-point support.}  Reading pour points from a shapefile
        enables batch catchment delineation in a single tool invocation.
  \item \textbf{Unified tool.}  Watershed and catchment are combined into one
        \texttt{WatershedWorkflowTool} rather than separate tools, reducing
        user friction and ensuring the flow direction/accumulation grids are
        computed only once.
\end{enumerate}


% ============================================================
%  ADDITIONAL CONTENT FROM Connections.md
%  ── Suggested placement: Section in Chapter 3 (Architecture)
%     or as a new section in Chapter 6 (Agent)
% ============================================================

% --- Paste the following into Chapter 3, Section "Data Flow" ---

% \section{Detailed Data Flow: Prompt to Result}
%
% The complete data flow from user input to rendered output traverses
% every major component of the architecture:
%
% \begin{enumerate}
%   \item User types in \texttt{prompt\_dock.py} \texttt{QLineEdit}.
%   \item \texttt{GeoSIDock.\_run\_query()} handles meta-commands
%         (\texttt{help}, \texttt{layers}, \texttt{tools}, \texttt{clear}).
%   \item \texttt{GeoSIDock.\_sync\_qgis\_layers()} copies
%         \texttt{QgsProject.mapLayers()} into the engine's state.
%   \item \texttt{\_QueryWorker} thread calls \texttt{GeoSI.query(text)}.
%   \item \texttt{agent.parse(text, state)} $\rightarrow$ \texttt{QueryParser}
%         $\rightarrow$ \texttt{AnalysisRequest}.
%   \item \texttt{agent.validate\_layers\_exist(req)} $\rightarrow$ returns
%         error if layer not found.
%   \item \texttt{agent.plan(req)} $\rightarrow$ \texttt{ExecutionPlan}
%         (Ollama-first).
%   \item \texttt{executor.run(plan)} $\rightarrow$ walks each
%         \texttt{ToolStep}, dispatching via \texttt{BackendTool} $\rightarrow$
%         \texttt{processing.run()}.
%   \item \texttt{ExecutionResult} returns to
%         \texttt{GeoSIDock.\_on\_result()}.
%   \item Chat dock renders answer, reasoning, and suggestions.
% \end{enumerate}


% --- Paste the following into Chapter 3 or Appendix ---

% \section{Layer Not Found Flow}
%
% When a user references a layer that is not loaded in the QGIS Layers panel,
% the engine responds with a clear error message listing available layers:
%
% \begin{lstlisting}[language={}]
% parser.parse("Buffer hospitals by 500m", ["Assets", "Kerala"])
%   -> request.entities["primary_layer"] = "hospitals"
%
% agent.validate_layers_exist(request, state)
%   -> primary "hospitals" not in available layers
%   -> find_similar_layers("hospitals", ["Assets","Kerala"]) == []
%   -> returns (False, "Layer 'hospitals' not found.
%                      Available layers: Assets, Kerala")
%
% agent.plan(...)   -> empty ExecutionPlan with error as reasoning
% executor.run(...) -> ExecutionResult(success=False, answer=error)
% prompt_dock       -> renders the error in red, lists loaded layers
% \end{lstlisting}
%
% The plugin \textbf{never} fetches data from OSM.  The user is always told
% to load the layer first.


% --- Paste the following into Chapter 6 (Agent) ---

% \section{Common Workflow Examples}
%
% \subsection{Workflow: Buffer + Intersect}
%
% \begin{lstlisting}[language={}]
% User: "Find residential zones within 500m of hospitals"
%
% parser:  intent=PROXIMITY,
%          entities={primary:hospitals, secondary:residential}
%          parameters={distance_value:500, distance_unit:meters}
% agent:   plan = [buffer(hospitals, 500),
%                  intersect(prev_output, residential)]
% exec:    buffer output registered as step_1_output;
%          intersect reads step_1_output + residential;
%          final Layer stored in state.
% \end{lstlisting}
%
% \subsection{Workflow: Raster NDVI}
%
% \begin{lstlisting}[language={}]
% User: "Calculate NDVI from Sentinel2_red and Sentinel2_nir"
%
% parser:  intent=RASTER, operation=ndvi
% agent:   plan = [ndvi(RED=Sentinel2_red, NIR=Sentinel2_nir)]
% exec:    BackendTool dispatches to "qgis" backend ->
%          processing.run("native:rastercalc", {...}) -> NDVI raster.
% \end{lstlisting}


% --- Paste into Appendix or Chapter 3 ---

% \section{Quick Reference: Key Classes and Methods}
%
% \begin{table}[ht]
% \centering
% \caption{GeoSI quick reference — key classes and entry points}
% \begin{tabular}{lll}
% \toprule
% \textbf{Need} & \textbf{Class / Module} & \textbf{Key Method} \\
% \midrule
% Run NL query          & \texttt{geosi\_engine.GeoSI}   & \texttt{.query(text)} \\
% List tools            & \texttt{ToolRegistry}          & \texttt{.list\_tools()} \\
% Run a tool directly   & \texttt{ToolRegistry}          & \texttt{.execute\_tool(name)} \\
% Parse text only       & \texttt{QueryParser}           & \texttt{.parse(text, layers)} \\
% Plan only             & \texttt{GeoSIAgent}            & \texttt{.plan(request)} \\
% Execute only          & \texttt{ExecutionEngine}       & \texttt{.run(plan)} \\
% Install QGIS backend  & \texttt{qgis\_bridge}          & \texttt{.install()} \\
% Chat dock UI          & \texttt{GeoSIDock}             & \texttt{\_run\_query()} \\
% \bottomrule
% \end{tabular}
% \end{table}


% --- Add to References ---

% \item O'Callaghan, J. F. and Mark, D. M., ``The Extraction of Drainage
% Networks from Digital Elevation Data,'' \textit{Computer Vision, Graphics,
% and Image Processing}, vol. 28, pp. 323--344, 1984.
%
% \item Jenson, S. K. and Domingue, J. O., ``Extracting Topographic Structure
% from Digital Elevation Data for Geographic Information System Analysis,''
% \textit{Photogrammetric Engineering and Remote Sensing}, vol. 54, no. 11,
% pp. 1593--1600, 1988.
%
% \item GDAL/OGR Contributors, ``GDAL/OGR Geospatial Data Abstraction
% Library,'' Open Source Geospatial Foundation, 2024.\\
% Available: \url{https://gdal.org}


% --- Add to Tool Registry Appendix ---

% watershed\_workflow & terrain & Native D8 watershed + catchment delineation \\
